\chapter{Justificación del modelo de desarrollo de software}
\label{Justificación del modelo de desarrollo de software}

\section{Modelo de desarrollo seleccionado}
El modelo de desarrollo seleccionado para este proyecto es Scrum, una metodología ágil que permite organizar el trabajo en ciclos cortos e iterativos denominados sprints. Esta elección resulta adecuada para el desarrollo de una plataforma web y aplicación móvil de gestión de fletes, debido a que el sistema cuenta con múltiples funcionalidades que pueden diseñarse, priorizarse y validarse de manera incremental.

Scrum permite dividir el proyecto en entregas parciales, facilitando que el equipo avance por módulos, tales como registro de usuarios, solicitud de fletes, comparación de transportistas, seguimiento GPS, chat interno, gestión de pagos, validación documental, panel administrativo y sistema de calificaciones. Esta forma de trabajo permite mantener un mejor control del avance y realizar ajustes a medida que se revisan los resultados del proyecto.

\section{Justificación de la elección}

La elección de Scrum se justifica principalmente por la naturaleza del software propuesto. Al tratarse de una plataforma con tres tipos de usuarios distintos —Cliente, Transportista y Administrador—, es necesario considerar distintas necesidades, flujos de interacción y funcionalidades. Esto hace conveniente utilizar un enfoque flexible, capaz de adaptarse a cambios o mejoras durante el desarrollo.

A diferencia de un modelo tradicional en cascada, donde las etapas se ejecutan de forma secuencial y con menor flexibilidad ante cambios, Scrum permite realizar revisiones frecuentes y recibir retroalimentación sobre el avance del producto. Esto es especialmente útil en proyectos donde la experiencia de usuario, el diseño de interfaces y la validación de funcionalidades son elementos importantes.

Además, Scrum favorece el trabajo colaborativo del grupo, ya que permite distribuir tareas, priorizar requerimientos y organizar entregas por períodos definidos. En este proyecto, dicha organización facilita que algunos integrantes trabajen en documentación, otros en diagramas, otros en mockups, otros en prototipo y otros en presentación, manteniendo una visión común del producto final.

\section{Aplicación de Scrum al proyecto}

Para aplicar Scrum en este proyecto se pueden considerar los siguientes roles:

\begin{itemize}
	\item \textbf{Product Owner:} David Eduardo Henriquez Bravo, integrante encargado de priorizar los requerimientos y representar las necesidades principales del usuario.
	\item \textbf{Scrum Master:} Lucas ``la cabra'' Contreras, integrante encargado de coordinar el avance del equipo, facilitar reuniones y remover obstáculos.
	\item \textbf{Equipo de desarrollo:} Todos los integrantes de este grupo, encargados de elaborar documentación, diagramas, mockups, prototipo, presentación y validaciones del sistema.
\end{itemize}

El trabajo puede organizarse mediante un Product Backlog, que corresponde a la lista priorizada de funcionalidades y tareas del proyecto. Algunas tareas del backlog son: registro de usuarios, solicitud de traslado, búsqueda de transportistas, seguimiento GPS, registro fotográfico, pagos, calificaciones, panel administrativo, diseño de mockups, diagramas y documentación del informe.

\section{Cronograma estimado del proyecto}

El cronograma estimado se organiza en sprints semanales, considerando las actividades necesarias para completar el informe, el prototipo dinámico y la presentación final.

\begin{table}[H]
	\centering 
	\caption{Tabla del cronograma del proyecto}
	\label{table: Tabla cronograma}
	\renewcommand{\arraystretch}{1.4}
	\begin{tabular}{|p{2cm}|p{10cm}|p{5cm}|}
		\hline 
		\centering\textbf{Sprint} & \centering\textbf{Actividad principal} & \centering\textbf{Responsable sugerido}\tabularnewline
		\hline 
		Sprint 1 & Definición del problema, usuarios, alcance y especificaciones generales de la aplicación. & Grupo completo \\
		\hline
		Sprint 2 & Levantamiento y redacción de requerimientos de usuario, funcionales y no funcionales. & Encargados de requerimientos \\
		\hline 
		Sprint 3 & Elaboración de diagramas BPMN, caso de uso principal y descripción de casos de uso. & Sepu, Lalo, Javier y Benjamin \\
		\hline 
		Sprint 4 & Diseño de mockups y prototipo dinámico de la plataforma. & Cony, Andy y Mati \\
		\hline 
		Sprint 5 & Desarrollo o armado del prototipo funcional, integración de pantallas y revisión de flujos & H, Dani y apoyo del grupo \\
		\hline 
		Sprint 6 & Redacción de secciones de soporte, seguridad, calidad, herramientas CASE y modelo de desarrollo. & Responsables por sección \\
		\hline 
		Sprint 7 & Revisión final del informe, referencias, formato, ortografía y consistencia. & Grupo completo\\
		\hline 
		Sprint 8 & Preparación de presentación, video/demo y ensayo del pitch final. & Sepu, H y grupo completo \\
		\hline 
	\end{tabular}
\end{table}




